%!TEX root = main.tex
%Observe that 

\subsection{Computational aspects I: Regularity with respect to \texorpdfstring{$\lambda$}{lambda}} The computational strategy will be to find the value $c_\lambda$ for a discrete subset of the possible $\lambda$, and prove regularity properties of $c_\lambda$ that guarantee that $C_{opt}= c_{\lambda^*}$ is not far from the maximum of the values of $c_{\lambda}$.

\begin{lemma}
The following properties hold:
\begin{itemize}
    \item[(i)] The Lipstchitz constant of $c_{\lambda}$ is at most $1$.
    \item[(ii)] If $C_{opt}=c_{\lambda^*}$, then
    $$
    c_{\lambda} \ge \frac{2 c_{\lambda^*} }{\lambda^{-1}\lambda^*+\lambda{\lambda^*}^{-1}}
    $$
\end{itemize}
\end{lemma}
\begin{proof}
%\newb{Clarify the proof of the first part of the lemma}
{\it{Proof of (i):}} Let $g$ be a function such that $\|g\|_1 = 1$, then $\langle g|K_w|g\rangle\le 1$, and 
$$
\left|\frac d {d\lambda} (\lambda \|g\|^2_2+ \lambda^{-1} \|g\|^2_1)^{-1}\right| = \lambda^ {-1}(\lambda \|g\|^2_2+ \lambda^{-1} \|g\|^2_1)^{-2} |\lambda \|g\|^2_2- \lambda^{-1}\|g\|^2_1|\le 1.
$$
The result follows from this.\\

{\it{Proof of (ii)}:} %For the second, 
We start observing that
$$
\lambda^{*}=\frac{\|f^*\|_1}{\|f^*\|_2},
$$
where $f^*$ is an extremizer for our original problem (This is when the equality happen in AM-GM inequality).
Then
\begin{align*}
 c_{\lambda}({\lambda^{*}}^{-1}\|f^*\|^2_1+\lambda^{*}\|f^*\|^2_{2})\geq 2\langle f^*|K_w|f^*\rangle= 2\|f^*\|_1\|f^*\|_2c_{\lambda^*}.
\end{align*}
Therefore
\begin{align*}
    c_{\lambda}\geq \frac{2\|f^*\|_1\|f^*\|_2c_{\lambda^*}}{{\lambda^{*}}^{-1}\|f^*\|^2_1+\lambda^{*}\|f^*\|^2_{2}}=\frac{2c_{\lambda^*} }{\lambda^{-1}\lambda^*+\lambda{\lambda^*}^{-1}}.
\end{align*}
\end{proof}


%\textbf{------------------}
\subsection{Computational aspects II: Discretizing the convolution kernel.}
Let $a,b\in \delta\mathbb{Z}$, and
as previously, let $V_\delta$ be the set of functions $g:[a,b)\to\mathbb{R}$ that are constants on intervals of length $\delta$.
Let $f:\delta \mathbb Z\cap [a,b)\to\mathbb{R}$, we define $\|f\|_{l^p_\delta}:= \left (\sum_{i\in\delta\mathbb Z\cap[a,b)} f(i)^p \delta\right )^{1/p}$ (essentially a $\delta-$discretization of the $L^p$ norm), and, we define $E[f]:=\sum_{i\in\delta\mathbb Z} f(i)\chi_{[i,i+\delta]}\in V_{\delta}$. Observe that we have the equality $\|f\|_{l^p_\delta} = \|E[f]\|_{L^p}$. For any $a,b\in \delta \mathbb Z$ this defines a natural isomorphism between $V_{\delta}$ and functions with domain $\delta \mathbb Z \cap [a,b)$. Let
$$
\iota :l^2(\delta \mathbb Z \cap [a,b)) \to L^2(\mathbb R \cap [a,b))\cap V_{\delta}
$$ be the map given by this isomorphism i.e $if=E[f]$ for all $f\in \delta\mathbb{Z}\cap[a,b)$ 
(Here we are considering $\delta \mathbb Z$ with $\delta$ times the counting measure as a measure). This isomorphism allows us to define the metrics induced in $V_{\delta}$ by  %$L^{1:2}$, 
$H_\lambda$, $B_\lambda$ as a metric for functions in $\delta \mathbb Z \cap [a,b)$. %Let $$\iota :l^2(\delta \mathbb Z \cap [a,b)) \to L^2(\mathbb R \cap [a,b))\cap V_{\delta}$$ be the map given by this isomorphism.


Given a function $w \in L^1 \cap L^{\infty} (
[a,b])$ there is a function in $\tilde w \in l^1 \cap l^{\infty} (\delta \mathbb Z \cap [a,b])$ such that for $f,g \in l^2(\delta \mathbb Z \cap [a,b))$ it holds that $$\langle  \iota f, w\ast_{\mathbb R} \iota g \rangle_{L^2(\mathbb R)} =  \langle  f, \tilde w\ast_{\delta \mathbb Z}  g \rangle_{l^2(\delta \mathbb Z)}. $$

The function $\tilde w$ is given explicitly by:

\begin{align*}
\tilde w(s) 
=&
\langle \delta^{-1} 1_s, \tilde w\ast_{\delta \mathbb Z} (\delta^{-1} 1_{0}) \rangle_{l^2(\delta \mathbb Z)} 
\\=&
\langle \delta^{-1} \chi_{[s,s+\delta)},  w\ast_{ \mathbb R} (\delta^{-1} \chi_{[0,\delta)}) \rangle_{L^2(\mathbb R)} 
%\\=&
%\langle \delta^{-1} 1_s, \tilde w\ast_{\delta \mathbb Z} (\delta^{-1} 1_{0}) \rangle_{l^2(\delta \mathbb Z)} 
\\=&
\fint_{s}^{s+\delta}\fint_{0}^{\delta} w(y-x)  dx dy
\\=&
\delta^{-2} \int_{s}^{s+\delta} w(t) (\delta-|t-s|) dt.
\end{align*}
In the two cases of special interest (when $w$ is a Gaussian function or when $w$ is the characteristic function of a set, we get more explicit values). In practice, however, the high stability of the integrals make it more accurate to perform the integrals numerically if $\delta$ is small than to compute the difference numerically (of the order of $10^{-5}$), for more details see the annotated code.



\subsection{Computational aspects II: Solving the problem for a fixed \texorpdfstring{$\lambda$}{lambda} at a discretization scale \texorpdfstring{$\delta$}{delta} }

What remains to do now is to give an algorithm that allows us to solve the extended problem for a fixed value of $\lambda$. The key fact that we use is that, if we %were able to 
remove the constraint $f\ge 0$, then the solution could be readily found using the power method for self-adjoint finite dimensional linear operators (symmetric matrices). We focus first in the situation when the constraint $f\ge 0$ is removed, and then argue that we can assume we are in that situation. \\

Recalling that
$$
C_{opt}(w) = \max_{\lambda>0} \max_{f \in L^1 \cap L^2([-a,a])}  2 \frac{\bra f w \ket f}{\lambda  \braket f f + \lambda^ {-1}\braket f 1 \braket 1 f}. 
$$
Assume that the functions under consideration are supported in $[-a,a)$, so we have $\braket 11 = 2a$. In an abuse of notation, we denote by $w$ the convolution operator associated to $w$.
The key observation is the following
$$
{\lambda \braket f f + \lambda^{-1} \braket f 1 \braket 1 f} = \bra f  (\sqrt \lambda Id + b_{\lambda} \ket 1 \bra 1 )^2 \ket f = \bra f A_{\lambda}^2 \ket f,
$$
where $b_{\lambda}$ is the unique positive solution to $\lambda^{-1} = 2\sqrt{\lambda} b_{\lambda}+2a b_{\lambda}^2$ and $A_{\lambda}$ is a hermitian positive definite matrix. Let $g:= A_{\lambda}f$, then we have:

$$
2 \frac{\bra f w \ket f}{\lambda  \braket f f + \lambda^ {-1}\braket f 1 \braket 1 f}  = 2 \frac{\bra g A_{\lambda}^{-1}  w A_{\lambda}^{-1}  \ket g}{\braket gg}
$$
defining $M_{\lambda}:= 2 A_{\lambda}^{-1}  w A_{\lambda}^{-1} $ we have that

$$
C_{opt}(w) = \max_{\lambda>0} \max \operatorname{Spec}(M_{\lambda}).
$$


%\subsection{Solving the auxiliary problem}

What remains to be shown is that the constraint $f\ge 0$ can indeed be dropped out. Let $f^*$ be the solution to the discretized problem constrained to $f\ge 0$. We can assume (by the discrete version of Riesz rearrangement inequality, see \cite[Chapter X]{HLP}) that $f^*$ is symmetric and non-increasing. Assume that, out of all potential extremizers to the discretized problem, $f^*$ has minimal support. \\

Let $\tilde V$ be the (finite dimensional) space of functions in $\delta \mathbb Z$ with the same support as $f^*$. Since $f^*$ extremizes the rayleigh quotient $\langle f|w|f\rangle$  it must be an eigenvector $w$ restricted to the support of $f$. It is the unique eigenvector: If there was another eigenvector $g^*$, the function $g^*-\alpha f^*$ would be an eigenvector as well, with a strictly smaller support if $\alpha$ is chosen appropriately.\\

Let $l$ be the number of $\delta$-intervals in the discretized shortest support. We can further constrain the optimization problem to functions supported in these $l$ intervals without changing the value of the optimization. In this case, the extremizer function $f^*$ will be an extremizer on the interior, and therefore, the largest eigenvalue of the bilinear form induced by $w$. We can find such eigenvalue by the power method.\\



This shows the following algorithm will give the extreme value up to discretization erros:

\begin{enumerate}
	\item Choose $\delta$ small enough for the desired error in Proposition \ref{bounding the error} to hold. Let $N$ be the number of $\delta-$intervals intersecting $[-a,a]$
	\item Fix $\lambda$ (and then repeat for a large enough set of $\lambda$ for the desired tolerance in Proposition \ref{bounding the error} to hold).
    \item For each natural $1\le k\le N$ solve the unconstrained discretized problem (the eigenvalue problem) with the power method. If the solution satisfies the constraints, keep as a potential problem.
    \item The maximum for our original problem (up to the tolerance/error arising from the discretization) is the maximum over all the solutions given in the previous step, where the maximum is taken over all the $\lambda$ and $k$.
\end{enumerate}

\subsection{A conjectured iterative method}


The Euler-Lagrange equations for $f$ in the support of $f$ can be written as:

	\begin{equation}\label{euler-lagrange-2}
		 \frac \optf {\|\optf\|_2^2} =  \max\left (\frac{2 \optf\ast w}{\int_{\mathbb R^2} \optf(x)\optf(y) w(x-y) dx dy}  - \frac 1 {\|\optf\|_1},0\right ).
	\end{equation}
This suggests a fixed point method (both in the discrete and continuous set-ups) to find an extremizer to the autoconvolution problem. This method converges in practice (see the last column in Table \ref{table}), and is hundreds of times faster than the method described in the previous section. We have however been unable to show convergence of such method. 

This method bears resemblance to the numerical method in \cite{CDGJLM}, where the authors were not able to show convergence of the method either.

\section{A comparison with the maximum problem}

At first glance this problem bears resemblance to the original maximum problem originally studied by Cilleruelo, Rusza and Vinuesa \cite{CRV}. The problem of the mean (when $w = \chi_{[-1/2,1/2]}$) was indeed posed by Steinberger as a relaxation of the problem of the maximum, namely finding the largest constant such that:

\begin{equation}
    \label{eq:maximo}
\text{max}_{-1/2 \le t \le 1/2} \int_{\R} f(t-x)f(x) \, dx \ge c_{\max} \left(\int_{-1/4}^{1/4} f(x) \, dx\right)^2.
\end{equation}

this problem was studied using computational methods by Cloninger and Steinberger \cite{CS}. The complexity of numerical method proposed in \cite{CS} grows exponentially in the level of discretization of the function, as opposed to the polynomial cost in the method proposed in this work.

One may hope that the methods in this work could generalize to methods in the maximum problem, however we believe that the problems are genuinely different for two reasons:

\begin{itemize}
    \item There is no reason to expect smooth (or Liptschitz) global extremizers to the maximum problem. In fact, the best known numerical extremizers \cite[Figure 1]{MV} seem non-smooth, and are certainly not symetrically decreasing. This prevents the gain of a $\delta^2$ in the discretization of the problem, which now has to be done much more carefully (\cite[Lemma 1]{CS}). 
    \item The (discretized) problem of the maximum seems to have multiple local extremizers of the functional in equation \eqref{eq:maximo}. This quite likely prevents fixed point methods to be efficient.  
\end{itemize}